% capitulos/cap01_por_que_programar.tex
\chapter{Por qué programar el método de rigidez}\label{cap:por_que_programar}

Empecemos por el final. El gráfico siguiente muestra el análisis completo de
un puente peatonal de 18 metros: arriba, cómo se deforma bajo su carga;
abajo, la fuerza axial de cada una de sus 21 barras, en azul las que trabajan
en tracción y en naranja las que trabajan en compresión.

\grafico{c12_analisis_c}

No lo dibujó un programa comercial. Lo dibujó una calculadora de unas
doscientas líneas de Python, escritas una por una a lo largo de este libro.
Al terminarlo, esa calculadora será la nuestra: la habremos escrito, entendido
y verificado paso a paso, y podremos usarla para analizar cualquier armadura
plana.

\begin{objetivos}
  \item Reconocer los límites del cálculo a mano de armaduras.
  \item Explicar qué hace por dentro un programa de análisis estructural.
  \item Describir los pasos del método de rigidez y su orden.
  \item Conocer los tres ejemplos que recorren el libro.
\end{objetivos}

% ---------------------------------------------------------------------------
\section{Los límites del cálculo a mano}\label{sec:limites}\index{método de los nudos}\index{isostática}

En los cursos de estática aprendimos a analizar armaduras con el método de
los nudos y el método de las secciones. Son métodos elegantes y conviene
dominarlos: con ellos verificaremos nuestra calculadora a lo largo del libro.
Pero tienen dos límites.

El primero es el \emph{tamaño}. El puente del gráfico tiene 12 nodos; el
equilibrio de cada nodo da dos ecuaciones, de modo que hay 24 ecuaciones con
24 incógnitas: 21 fuerzas en las barras y 3 reacciones. Resolverlas a mano es
posible, pero largo y propenso a errores, y un error en un nudo se arrastra a
todos los siguientes. Una armadura de cubierta de una nave industrial puede
tener cientos de barras.

El segundo límite es más profundo. Una armadura de $b$ barras, $r$
reacciones y $n$ nodos tiene $b + r$ fuerzas desconocidas y solo $2n$
ecuaciones de equilibrio. Si $b + r > 2n$, la armadura es
\termino{hiperestática}{statically indeterminate}: el equilibrio por sí solo
no basta para obtener sus fuerzas, porque estas dependen además de la
rigidez de las barras. Es el caso de un cuadro con sus dos diagonales, que
estudiaremos como ejemplo B: 6 barras y 3 reacciones frente a 8 ecuaciones.
Ningún método basado solo en el equilibrio puede resolverlo.

El \termino{método de rigidez}{stiffness method} supera ambos límites. Toma
como incógnitas los desplazamientos de los nodos en lugar de las fuerzas,
relaciona cada desplazamiento con las fuerzas que produce mediante la rigidez
de las barras, y reúne todo en un sistema de ecuaciones lineales,
\begin{ecuacionclave}
  \mK\,\vU = \vF,
\end{ecuacionclave}
donde $\mK$ es la \termino{matriz de rigidez}{stiffness matrix} de la
armadura, $\vU$ el vector de desplazamientos de los nodos y $\vF$ el de las
fuerzas aplicadas sobre ellos. El método es el mismo para una armadura
isostática o hiperestática, de tres barras o de trescientas. Y, sobre todo,
es sistemático: cada paso se repite igual para cada barra y cada nodo, lo
que lo hace ideal para un computador.

% ---------------------------------------------------------------------------
\section{Qué hay dentro de un programa comercial}\label{sec:programas}

Los programas de análisis estructural que se usan en la práctica, como
SAP2000, ETABS o Ftool, resuelven las estructuras con el método de rigidez,
generalizado a barras que además se flexionan. Cuando dibujamos una armadura
en uno de ellos, el programa hace lo mismo que haremos en este libro: calcula
la rigidez de cada barra, ensambla la matriz de la estructura, resuelve el
sistema y obtiene las fuerzas de cada barra.

Usar un programa sin entender lo que hace por dentro es riesgoso. El programa
siempre da un resultado, también cuando los datos están mal: un apoyo
olvidado, una carga en la unidad equivocada o un nodo desplazado. Reconocer
que un resultado no tiene sentido exige saber de dónde sale. Por eso este
libro no enseña a usar un programa, sino a construirlo: al escribir cada paso
del método, sabremos qué esperar de él y cómo verificarlo.

\begin{parasabermas}
El método de rigidez se desarrolló en la industria aeronáutica en la década
de 1950, junto con los primeros computadores, y está en la base del método
de los elementos finitos. Su desarrollo completo para armaduras, vigas y
pórticos está en \textcite{kassimali2012} y en \textcite{mcguire2000}.
\end{parasabermas}

% ---------------------------------------------------------------------------
\section{El mapa del método}\label{sec:mapa}\index{método de rigidez!pasos del}

La \cref{fig:mapa_metodo} resume el método de rigidez para armaduras planas
y es, a la vez, el mapa de este libro: cada paso es una parte de nuestra
calculadora, y cada capítulo construye uno o dos pasos.

\begin{figure}[ht]
  \centering
  \figura{mapa_metodo}
  \caption{Los pasos del método de rigidez, con la ecuación principal de
    cada uno y el capítulo donde lo programamos.}
  \label{fig:mapa_metodo}
\end{figure}

En palabras:
\begin{enumerate}
  \item Describimos la armadura con datos: las coordenadas de sus nodos, las
    barras que los unen, los apoyos y las cargas.
  \item Calculamos la longitud y la orientación de cada barra.
  \item Con ellas obtenemos la matriz de rigidez de cada barra, que relaciona
    los desplazamientos de sus extremos con las fuerzas en ellos.
  \item Numeramos los desplazamientos posibles de los nodos, sus
    \termino{grados de libertad}{degrees of freedom}.
  \item Ensamblamos las matrices de las barras en la matriz de la armadura, y
    las cargas en el vector de fuerzas.
  \item Separamos los grados de libertad libres de los restringidos por los
    apoyos.
  \item Resolvemos el sistema: obtenemos los desplazamientos de los nodos y,
    con ellos, las reacciones.
  \item De los desplazamientos obtenemos la fuerza axial de cada barra.
  \item Presentamos, verificamos y dibujamos los resultados.
\end{enumerate}

\begin{recuerda}
Todo el método se basa en tres principios que ya conocemos: el
\emph{equilibrio} de cada nodo, la \emph{compatibilidad} (las barras unidas
en un nodo se desplazan con él) y la \emph{ley de Hooke} (cada barra se
alarga en proporción a su fuerza). Supondremos, como en resistencia de
materiales, que el material es elástico lineal y que los desplazamientos son
pequeños.
\end{recuerda}

% ---------------------------------------------------------------------------
\section{Tres armaduras para aprender}\label{sec:tres_ejemplos}\index{ejemplo A}\index{ejemplo B}\index{ejemplo C}

A lo largo del libro usaremos tres armaduras de acero, con datos realistas
($E = \SI{200}{\giga\pascal}$), que se resumen en la
\cref{tab:tres_ejemplos}. La columna $g$ es el grado de indeterminación,
$g = b + r - 2n$: es cero en una armadura isostática y positivo en una
hiperestática.

\begin{table}[ht]
  \centering
  \caption{Los ejemplos del libro.}
  \label{tab:tres_ejemplos}
  \small
  \begin{tabular}{@{}lcccl@{}}
    \toprule
    Ejemplo & Nodos & Barras & $g$ & Para qué sirve \\
    \midrule
    A. Triángulo & 3 & 3 & 0 & Verificar cada paso a mano \\
    B. Cuadro arriostrado & 4 & 6 & 1 & Un caso hiperestático \\
    C. Puente Pratt & 12 & 21 & 0 & Una armadura de verdad \\
    \bottomrule
  \end{tabular}
\end{table}

El \emph{ejemplo A} es un triángulo de \SI{6}{\metre} de base y
\SI{4}{\metre} de altura, lo bastante pequeño para comprobar a mano cada
número que produzca la calculadora. El \emph{ejemplo B} es el cuadro con dos
diagonales del que hablamos antes: nos recordará que el método resuelve lo
que el equilibrio por sí solo no puede. Y el \emph{ejemplo C} es el puente
del comienzo del capítulo, la prueba final de la calculadora.

Que haya tres ejemplos distintos no es un capricho. Un programa puede dar el
resultado correcto con un ejemplo por casualidad, por ejemplo si confunde dos
números que en ese ejemplo resultan iguales. Probarlo con armaduras de
tamaños y formas distintas hace que esos errores salgan a la luz.

% ---------------------------------------------------------------------------
\section{Lo que necesitamos para empezar}\label{sec:requisitos}

Este libro supone que conocemos la estática de armaduras, las nociones de
esfuerzo y deformación de resistencia de materiales, y el álgebra de
matrices: producto, transpuesta e inversa. No supone ningún conocimiento de
programación. Aprenderemos Python a medida que lo necesitemos, siempre con
ejemplos de armaduras, y cada herramienta nueva aparecerá en una caja como
esta:

\begin{python}
Las cajas de Python explican las herramientas del lenguaje que usamos por
primera vez. El anexo~A las reúne todas para consultarlas rápidamente.
\end{python}

Tampoco necesitamos instalar nada: escribiremos nuestro código en notebooks de
Google Colab, que funcionan en el navegador. En el \cref{cap:primer_notebook}
creamos el primero.

% ---------------------------------------------------------------------------
\repasoyresumen

\begin{itemize}
  \item El cálculo a mano es largo para armaduras grandes e insuficiente para
    las hiperestáticas, cuyo grado de indeterminación $g = b + r - 2n$ es
    positivo.
  \item El método de rigidez toma como incógnitas los desplazamientos de los
    nodos y resuelve el sistema $\mK\,\vU = \vF$. Es el mismo para cualquier
    armadura y es el que usan los programas comerciales.
  \item El método tiene nueve pasos (\cref{fig:mapa_metodo}), que
    programaremos en los capítulos 2 a 12.
  \item Los ejemplos A, B y C, de tamaños y propiedades distintos, nos
    permitirán verificar la calculadora.
\end{itemize}

% ---------------------------------------------------------------------------
\ejercicios

\begin{ejercicio}{1.1}
Calcular el grado de indeterminación de los ejemplos A, B y C, sabiendo que
los tres tienen un apoyo articulado y uno de rodillo. ¿Cuáles pueden
resolverse con el método de los nudos?
\end{ejercicio}

\begin{ejercicio}{1.2}
Cada nodo de una armadura plana puede desplazarse en dos direcciones.
¿Cuántas filas y columnas tiene la matriz $\mK$ de cada ejemplo? ¿Cuántos
números tiene la del ejemplo C?
\end{ejercicio}

\begin{ejercicio}{1.3}
El puente del ejemplo C tiene seis paneles. Para un puente del mismo tipo
con $p$ paneles, deducir cuántos nodos y barras tiene, y comprobar que sigue
siendo isostático. ¿De qué tamaño sería la matriz $\mK$ de un puente de 20
paneles?
\end{ejercicio}
